\chapter{คลังรหัสและกฎเกณฑ์สกิลการวิจัยวิชาการ การสังเคราะห์ และการเขียน}

ภาคผนวกนี้รวบรวมข้อมูลจำเพาะ กฎเกณฑ์อัลกอริทึม และพรอมต์สถาปัตยกรรมของสกิลด้านการวิจัย การทบทวนวรรณกรรม และการขัดเกลาสำนวนวิชาการประจำเล่มที่ 2

\section{สกิล: anti-ai-slop}

\begin{teachingplan}{ข้อมูลจำเพาะและเมทาดาทาของสกิล}
\textbf{ชื่อสกิล:} \texttt{anti-ai-slop}\\
\textbf{วัตถุประสงค์:} ตรวจจับและขจัดคำฟุ่มเฟือยของ AI ปรับจังหวะประโยค และรักษาน้ำเสียงมนุษย์\\
\textbf{กลไกควบคุม:} Vocabulary Denylist, Burstiness Calibration, Warranted Uncertainty
\end{teachingplan}

\subsection{กฎเกณฑ์บัญชีดำคำศัพท์ AI (Denylist Rules)}

\begin{lstlisting}[language=bash, caption={Anti-AI Slop Denylist Dictionary Configuration}]
DENYLIST_WORDS = [
  "delve", "tapestry", "crucial", "pivotal", "seamless", 
  "robust", "leverage", "elevate", "navigate", "testament",
  "beacon", "transformative", "game-changer", "foster"
]

SYMMETRIC_PATTERNS = [
  "Not only X, but also Y",
  "It is important to note that",
  "In today's fast-paced digital world"
]
\end{lstlisting}

\section{สกิล: research-synthesis}

\begin{teachingplan}{ข้อมูลจำเพาะและเมทาดาทาของสกิล}
\textbf{ชื่อสกิล:} \texttt{research-synthesis}\\
\textbf{วัตถุประสงค์:} สังเคราะห์หลักฐานงานวิจัยระดับ Peer-Review ด้วยตาราง 7 คอลัมน์\\
\textbf{มิติการวิเคราะห์:} Claim, Method, Evidence, Limitations, Authority, Recency, Network
\end{teachingplan}

\subsection{โครงสร้างชุดข้อมูลเมทริกซ์ 7 คอลัมน์ (JSON Schema)}

\begin{lstlisting}[language=JavaScript, caption={7-Column Synthesis Matrix Schema}]
const ResearchMatrixSchema = {
  paper_id: "string",
  authors: "array",
  year: "number",
  claim: "string",        // Primary asserted finding
  method: "string",       // Experimental setup & sample size
  evidence: "object",     // Metrics: p-value, effect size, std
  limitations: "array",   // Uncontrolled variables & biases
  authority: "string",    // Author citation index & affiliation
  recency_score: "number",// 1 to 5 scale
  citation_network: "array" // Key foundational papers cited
};
\end{lstlisting}

\section{สกิล: citation-formatter}

\begin{teachingplan}{ข้อมูลจำเพาะและเมทาดาทาของสกิล}
\textbf{ชื่อสกิล:} \texttt{citation-formatter}\\
\textbf{วัตถุประสงค์:} จัดรูปแบบบรรณานุกรมอัตโนมัติตามมาตรฐานสากล (APA 7th, IEEE, Harvard)\\
\textbf{ฟังก์ชันการทำงาน:} DOI Verification, Author Name Inversion, Journal Title Formatting
\end{teachingplan}

\subsection{อัลกอริทึมจัดรูปแบบ APA 7th Edition}

\begin{lstlisting}[language=JavaScript, caption={APA 7th Reference Generator}]
function formatAPA7(entry) {
  // entry: { authors: [{first, last}], year, title, journal, volume, issue, pages, doi }
  const authorStr = entry.authors.map(a => `${a.last}, ${a.first[0]}.`).join(", ");
  const issueStr = entry.issue ? `(${entry.issue})` : "";
  const doiStr = entry.doi ? ` https://doi.org/${entry.doi}` : "";
  
  return `${authorStr} (${entry.year}). ${entry.title}. ${entry.journal}, ${entry.volume}${issueStr}, ${entry.pages}.${doiStr}`;
}
\end{lstlisting}
